% !PRINTED 107
\weberpage{107}
Nehmen wir als Beispiel $m = 48$, so haben wir zwei Möglichkeiten, nämlich:
\[ \begin{aligned}
e_{1,1} &= 1, & e_{1,2} &= 0, & e_{1,3} &= 1\\
e_{1,1} &= 1, & e_{1,2} &= 1, & e_{1,3} &= 1,
\end{aligned} \]
und man findet aus ihnen die beiden Gruppen:
\[ \begin{aligned}
\mathfrak{A}_1 &= 1,\ 17,\ 31,\ 47\\
\mathfrak{A}_2 &= 1,\ \ 7,\ 41,\ 47.
\end{aligned} \]

\begin{center}
\S.\,25.
\end{center}

\begin{center}
\gesperrt{Cubische Abel'sche Gleichungen.}
\end{center}

Wir wollen diese Betrachtungen über die cubischen und
biquadratischen Kreistheilungsgleichungen mit dem Beweise eines
merkwürdigen Satzes abschliessen, der geeignet ist, diesen Gleichungen
 ein weit höheres und allgemeineres Interesse zu verleihen
 und der ein specieller Fall eines ganz allgemeinen Satzes
ist, den wir erst später kennen lernen werden.

Der Satz lautet so:

\gesperrt{Alle cubischen und biquadratischen Abel'schen
Gleichungen im Körper der rationalen Zahlen sind
Kreistheilungsgleichungen.}

Dadurch gewinnen die Resultate der beiden letzten Paragraphen
 einen höheren Werth. Es folgt dann nämlich, dass
durch die dort gebildeten Kreistheilungsperioden $\eta$ \gesperrt{alle Abel'schen
Zahlkörper} dritten und vierten Grades dargestellt sind,
dass andere als die dort besprochenen Körper dieser Art nicht
existiren.

Die Möglichkeit, diesen Satz für die speciellen Fälle der
cubischen und biquadratischen Gleichungen einfach und ohne
weitere Vorbereitung zu beweisen, beruht darauf, dass in den
Körpern der dritten und vierten Einheitswurzeln $R(\varrho)$, $R(i)$ dieselben
 Gesetze der Zerlegung der Zahlen in ihre Primfactoren
gelten, wie in den Körpern der rationalen Zahlen, wie in den
\S\S.\,173, 174 des ersten Bandes nachgewiesen ist. Es ist darum
auch von Interesse, diese besonderen Fälle genauer zu betrachten,
weil dadurch der Gang und das Ziel des später beizubringenden
 allgemeinen Beweises deutlicher erkannt wird.
